% 本文件由 LaTeX 排版 Agent 根据有效 translation.json 自动生成。
% author=Aristides the Philosopher; work_id=1012; title=阿里斯蒂德斯的辩护

\chapter{阿里斯蒂德斯的辩护}

% structure: p0001 source=h1 target=omitted reason=page-title-already-rendered-by-parent

\Emph{以下是哲学家\Person{阿里斯蒂德斯}在\Person{哈德良}王面前为尊崇天主所作的辩护。}\par

……全能的凯撒\Person{提图斯·哈德良·安东尼努斯}，可敬而仁慈的，来自雅典哲学家\Person{马尔西安·阿里斯蒂德斯}。\par

1. 国王啊，我因天主的\OriginalTerm{恩宠}{grace}来到这个世界；当我观察了天、地、海，并察看了太阳和其余受造之物时，我惊叹于世界的美丽。我认识到世界及其中的一切，都由另一位的权能所推动；我明白那推动它们的，就是\OriginalTerm{天主}{God}，衪隐藏在其中，并被它们所遮蔽。显然，那引起运动者比被动者更有能力。但是，要探究这万有的推动者，衪的本性如何（因为在我看来，衪的本性实在无法测度），并要论辩衪治理的不变恒常，以便完全把握——这对我而言是徒劳的；因为人不可能完全领悟。然而，关于这世界的推动者，我要说，衪是万有之天主，为了人类而造了万物。在我看来，这是合理的：人应当敬畏天主，不应欺压人。\par

因此我说，\OriginalTerm{天主}{God}不受生，也不受造，是恒存的本性，无始无终，不死不灭，完美的，不可理解。当我说衪是\Quote{完美的，}时，意思是在衪内没有任何缺陷，衪不需要任何事物，但万物都需要衪。当我说衪是\Quote{没有开端的，}时，意思是凡有开端的也必有终结，而有终结的可以归于终结。衪没有名字，因为一切有名字的都与受造之物同类。衪没有形式，也没有肢体的组合；因为凡拥有这些的，都与被塑造之物同类。衪既非男，也非女。诸天不限制衪，相反，诸天与一切有形的、无形的，都从衪获得界限。衪没有对手，因为不存在任何比衪更强的。衪没有愤怒和怨恨，因为没有什么是能抵挡衪的。无知和遗忘不在衪的本性内，因为衪全然是智慧和明悟；在衪内，一切存在得以稳固。衪不需要祭品和奠酒，甚至任何可见之物；衪不需要任何人的什么，而一切受造物都需要衪。\par

2. 既然我们已就天主向你们谈论了这么多，尽我们言词所能及，现在就让我们转向人类，好能知道他们中谁分享了我们所讲的真理，谁背离了它。\par

国王啊，你清楚知道，世上有四类人：蛮族、希腊人、犹太人和基督徒。蛮族追溯其宗教的起源，来自\Person{克洛诺斯}和\Person{瑞亚}及他们的其他神祇；希腊人则追溯自\Person{赫勒诺斯}，据说他出自\Person{宙斯}。由\Person{赫勒诺斯}而生\Person{埃俄洛斯}和\Person{克苏托斯}；另有些人则是\Person{伊那科斯}和\Person{福洛纽斯}的后裔，最后还有出自埃及的\Person{达纳奥斯}、出自\Person{卡德摩斯}以及\Person{狄俄尼索斯}。\par

犹太人则追溯其种族的起源，来自\Person{亚巴郎}，他生了\Person{依撒格}，\Person{依撒格}生了\Person{雅各伯}。雅各伯生了十二个儿子，他们从\Place{叙利亚}迁徙到\Place{埃及}；在那里，他们因那为他们制定法律的，被称为希伯来民族；最后他们被称为犹太人。\par

基督徒则追溯其宗教的开端，来自\Person{耶稣}\OriginalTerm{默西亚}{Messiah}；衪被称为至高者天主的儿子。据说，天主从天上降下，并由一位希伯来童贞女取得了肉躯，穿上了血肉之体；天主子居住在人女之中。这就是所谓的\OriginalTerm{福音}{gospel}所教导的，不久前在他们中间宣讲的；你如果阅读它，也能明认其中的能力。这位耶稣，出生自希伯来民族；衪有十二位门徒，为要使衪\OriginalTerm{降生成人}{incarnation}的旨意按时完成。但衪自己被犹太人刺透，死了并被埋葬；他们说，三天后衪复活并升了天。于是，这十二位门徒走遍了已知的世界各地，以完全的谦逊和正直持续彰显衪的伟大。因此，现今凡相信这宣讲的人，都称为基督徒，且已闻名于世。\par

所以，如上所述，有四类人：蛮族、希腊人、犹太人和基督徒。\par

此外，风服从天主，火服从天使；水服从魔鬼，地服从世人。\par

3. 那么，我们就从蛮族开始，然后依次论及各民族，好能看出他们中谁持有关于天主的真理，谁持有错误。\par

蛮族因为不认识天主，便在元素中迷失了，开始朝拜受造之物，而不敬拜造物主；为此，他们制造了偶像，将它们关在神龛中，看哪！他们朝拜它们，同时小心翼翼地守护着，免得它们的神被强盗偷去。蛮族没有意识到，那守护者大于被守护者，且凡创造者大于受造者。那么，如果他们的神祇过于软弱，连自身安全都保不住，又怎能顾及人的安全呢？所以，蛮族陷入了何等大的错误，朝拜那些无生命的、不能帮助他们丝毫的偶像。国王啊，我不禁对他们的哲学家感到惊奇，连他们也误入歧途，竟将那些为尊崇元素而造的偶像称为神；而他们的贤士竟没有领悟到，元素本身也是可分解、可朽坏的。因为如果某个元素的一小部分可分解或毁灭，那么整体也可分解和毁灭。那么，如果元素本身是可分解、可毁灭的，且被迫服从另一位比它们更强有力的，如果它们本性上并非神，那么他们竟然将那些为尊崇它们而造的偶像称为神，是何道理呢？所以，他们中间的哲学家给追随者带来的错误是何等大啊。\par

4. 王啊，现在让我们转向元素本身，以便我们指明它们并非神，而是受造之物，容易朽坏和变化，与人的本性相同；而天主则不灭不变，不可见，虽不可见却洞察、掌管并变化万物。\par

那些相信大地是神的人至今欺骗了自己，因为大地被犁耕、种植和挖掘；它承受人与牲畜、家畜的污秽。有时它变得不结果实，因为若被烧成灰烬，便失去生命，从瓦罐中不会发芽。此外，若水积聚其上，它便与出产一同瓦解。它被人和牲畜践踏，承接被杀者的血污；它被挖开，填满死人，成为尸体的坟墓。但是，一个圣洁、尊贵、受祝福而不朽的本性，不可能允许这些事中的任何一件。由此我们认为，大地不是神，而是天主的化工。\par

5. 同样地，那些相信水是神的人也误入歧途。因为水是为人的用途而受造，并在多方面受人支配。水会遭受变化、容纳不洁，在煮沸成多种物质时被毁坏并失去其本性。它们染上不属于自己的颜色；也被霜冻凝结，混杂并渗透着人与牲畜的污秽及被杀者的血。它们被巧匠用导水渠遏制，逆其性而流淌并分流，进入园圃和其他地方，为使人收集并涌出作为丰产的媒介，并洗净一切不洁，履行人对它们的需求。因此水不能是神，而是天主的化工及世界的部分。\par

同样地，那些相信火是神的人也犯了不小的错误。因为火也是为人服务而受造，并在多方面受人支配——烹煮食物、作为铸造金属的手段，以及其他陛下所知的用途。同时，火在一些方式下被熄灭和扑灭。\par

再者，那些相信风的运动是神的人也错了。因为我们清楚知道，这些风是在另一位的权下，其运动有时增强，有时因那管制它们的命令而减弱和停止。因为风是天主为了人而造，为供给树木、果实和种子的需要；为将运载人的必需品和货物从产地运到不产之地的海上船只带来；并治理世界的各方。至于风自身，它有时增强又减弱；在管制它们的那一位的命令下，在一处带来助益，在另一处造成灾祸。人类也能用已知的方法限制和控制它，以便它为人们提供所需的功用。风自身毫无权柄。因此风不可称为神，而是天主所造的物。\par

6. 同样，那些相信太阳是神的人也错了。因为我们看到它被另一位的强制力所驱动，旋转运行，从一宫到另一宫，每日升起落下，为植物树木的生长提供温暖，并把它所混合（阳光）的地上每一生长之物带入空中。在它的运行中，它与其余星辰相比有其共同的部分；尽管本性单一，它却与许多部分相连，以供应人的需要；而且这不是出于它自己的意志，而是出于统治它者的意志。因此太阳不能是神，而是天主的化工；同样，月亮和星辰也是如此。\par

7. 那些相信古人中有些是神的人，也是大错特错了。因为正如您自己承认的，王啊，人是由四大元素以及一个灵魂和一个精神构成的（因此他被称为\OriginalTerm{小宇宙}{microcosm}），缺少任何一个部分，他就不能成立。他有起始和终结，有生有死。但天主，如我所言，其本性并无这些事，而是非受造和不灭的。因此，我们不可能把人树立为具有天主本性——这人时而盼望着欢乐，却遭遇忧患；时而期盼欢笑，却迎来哭泣——他易怒、贪婪、嫉妒，还有其他缺陷。他在许多方式下被元素以及动物所毁灭。\par

因此，王啊，我们必须认清蛮族的错误，即因他们未曾寻见真天主的踪迹，便偏离真理，随从自己幻想的欲念，事奉可朽的元素和无生命的偶像，由于他们的错误，未能领悟真天主是谁。\par

8. 让我们也转向希腊人，好知道他们对真\OriginalTerm{天主}{God}持何种看法。希腊人因为比蛮族更精巧，就比蛮族偏离得更远；他们引入了许多虚构的神明，立了一些为男神，一些为女神；在他们的神中，有些被发现行奸淫、杀人、受迷惑、嫉妒、忿怒、纵欲、弑亲、偷盗、抢劫。他们说，有些神残废跛脚，有些行巫术，有些竟发了疯，有些弹奏里拉，有些喜爱在山间游荡，有些甚至死了，有些被闪电劈死，有些作了人的奴仆，有些逃跑，有些被人绑架，有些竟被人哀悼悲叹。他们说，有些下到\OriginalTerm{阴府}{Sheol}，有些受重伤，有些变为动物的模样以引诱凡间女子，有些与男性同寝玷污自己。他们说，有些与自己的母亲、姐妹和女儿成婚。他们说他们的神与人的女儿行奸淫；从这些生出了一族，也是会死的。他们说，有些女神为美貌起争执，出现在人前受评判。国王啊，希腊人就这样提出了关于他们的神和他们自己的污秽、荒谬和愚妄，竟称那些本不是神的本性如此的为神。因此，人类受到鼓励去行奸淫、邪淫、偷窃以及一切可憎、可恨、可恶的事。因为如果他们称为神的那些都行了上面所记的这一切事，人岂不是更要行这些事——那些相信他们的神自己曾行这些事的人。由于这一错误的污秽，人类遭遇了扰攘的战争、大饥荒、苦为掳囚及全然荒凉。看啊！他们受苦，这一切临到他们，都是因这缘故；他们虽忍受这些事，心里却不省悟这些事临到他们，是因他们的错谬。\par

9. 让我们进一步看他们关于他们神明的记载，以便仔细证明上面所说的一切。首先，希腊人提出一位神\OriginalTerm{克洛诺斯}{Kronos}，即\OriginalTerm{基芸}{Chiun}（\OriginalTerm{撒图尔努斯}{Saturn}）。敬拜他的人将儿女献祭给他，有的还为他活活焚烧，以表敬意。他们说他在其妻\OriginalTerm{瑞亚}{Rhea}之外，又从她生了众多子女。他也从她生了\OriginalTerm{狄奥斯}{Dios}，即被称为\OriginalTerm{宙斯}{Zeus}的。最后他（克洛诺斯）发了疯，因畏惧一个向他显示的谕示，便开始吞吃自己的儿子们。宙斯被偷偷从他身边带走，他竟不知道；最后宙斯捆住他，割掉他男性的记号，抛入海中。于是，如他们在神话中所说，就生出了\OriginalTerm{阿佛洛狄忒}{Aphrodite}，即被称为\OriginalTerm{阿斯塔蒂}{Astarte}的。他（宙斯）将克洛诺斯上了锁链，丢到黑暗中。国王啊，希腊人关于他们众神之首所提出的，是何等大的错谬和耻辱，他们竟说到他这一切。神是不可能被捆绑或残缺的；若非如此，那神就真是可怜的了。\par

继克洛诺斯之后，他们又提出另一位神宙斯。他们说他掌握了主权，成为众神之王。他们说他变化为兽形及其他形状，为引诱凡间女子，并从她们为自己生子。他们说，有一次他因爱\OriginalTerm{欧罗巴}{Europe}和\OriginalTerm{帕西淮}{Pasiphae}而变成公牛，又因爱\OriginalTerm{达那厄}{Danae}而变成金子，因爱\OriginalTerm{勒达}{Leda}而变成天鹅，因爱\OriginalTerm{安提俄珀}{Antiope}而变成男子，因爱\OriginalTerm{路娜}{Luna}而变成闪电，他由此生了众多儿女。他们说，他由安提俄珀生了\OriginalTerm{泽托斯}{Zethus}和\OriginalTerm{安菲翁}{Amphion}，由路娜生了\OriginalTerm{狄俄尼索斯}{Dionysos}，由\OriginalTerm{阿尔克墨涅}{Alcmena}生了\OriginalTerm{赫拉克勒斯}{Hercules}，由\OriginalTerm{勒托}{Leto}生了\OriginalTerm{阿波罗}{Apollo}和\OriginalTerm{阿耳忒弥斯}{Artemis}，由达那厄生了\OriginalTerm{珀耳修斯}{Perseus}，由勒达生了\OriginalTerm{卡斯托耳}{Castor}、\OriginalTerm{波吕丢刻斯}{Polydeuces}、\OriginalTerm{海伦}{Helene}和\OriginalTerm{帕鲁杜斯}{Paludus}，由\OriginalTerm{谟涅摩绪涅}{Mnemosyne}生了九个女儿，她们被称为\OriginalTerm{缪斯}{Muses}，由欧罗巴生了\OriginalTerm{米诺斯}{Minos}、\OriginalTerm{拉达曼托斯}{Rhadamanthos}和\OriginalTerm{萨耳珀冬}{Sarpedon}。最后，他因贪恋牧人\OriginalTerm{伽倪墨得斯}{Ganydemos}，变成了一只鹰的模样。\par

国王啊，因着这些故事，人间生出了许多邪恶，直到今日，人仍在效法他们的神，行奸淫，与母亲、姐妹同寝玷污自己，与男性同寝，还有些竟敢杀害父母。因为如果那被称作众神之首和王的尚且行这些事，敬拜他的人岂不更要效法他呢？希腊人在关于他的叙述中提出了极大的愚妄。因为神是不可能行奸淫、邪淫、亲近男性或杀害父母的；若非如此，他就比一个行毁灭的魔鬼更不如了。\par

10. 他们又提出另一位神\OriginalTerm{赫淮斯托斯}{Hephaistos}。他们说他瘸腿，头上戴着帽子，手持火钳和锤子；他从事打铁的工艺，以便谋生所需。那么这位神竟是如此贫乏吗？但神不可能贫乏或瘸腿，否则他就毫无价值了。\par

他们又引进另一位神，称为\OriginalTerm{赫耳墨斯}{Hermes}。他们说他是个贼，贪恋吝啬，贪图利益，又是个术士，残缺不全，还是个运动员，又是通译语言的。但神不可能是术士、贪吝、残废、贪求非己之物或作运动员的。若非如此，他就显得无用了。\par

在他们之后，他们又提出另一位神，名叫\Person{阿斯克勒庇俄斯}。他们说他是一位医师，制备药物和膏药，以供应生活所需。那么这位神是有所匮乏吗？后来，他因\Place{拉刻代蒙}的\Person{廷达瑞俄斯}而被\Person{宙斯}用雷电击打，就这样死了。如果\Person{阿斯克勒庇俄斯}是神，当他被雷击时都无法自救，又如何能帮助别人呢？但一个神性竟会匮乏或被雷电摧毁，这是不可能的。\par

他们又提出另一位作神，称他为\Person{阿瑞斯}。他们说他是一名战士，嫉妒，贪求绵羊和本不属于他的东西。他以武力谋利。据说最后他与\Person{阿佛洛狄忒}通奸，被小孩子\Person{厄洛斯}和\Person{阿佛洛狄忒}的丈夫\Person{赫淮斯托斯}捉住。但神不可能是一名战士、受捆绑或犯奸淫。\par

他们又说\Person{狄俄倪索斯}，哼！是神，他安排夜间的狂欢宴饮，教人醉酒，抢夺属于别人的妇女。最后，他们说，他发了疯，遣散使女，逃入旷野；在疯狂中他吃蛇。最终他被\Person{提坦}杀死。如果\Person{狄俄倪索斯}是神，当他被杀时不能自救，又怎能帮助别人呢？\par

接下来他们提出\Person{赫拉克勒斯}，说他是神，他憎恶可憎之物，是暴君、战士、瘟疫的毁灭者。他们也说他最后发了疯，杀死自己的孩子，投身火中而死。如果\Person{赫拉克勒斯}是神，在所有这些灾难中都无法自救，别人怎能向他求助呢？但神不可能发疯、醉酒、杀害子女，或被火焚毁。\par

11. 在他们之后，他们又提出另一位神，称他为\Person{阿波罗}。他们说他嫉妒且反复无常，有时手持弓箭和箭囊，有时又拿着竖琴和琴拨。他向人发布神谕，为要从他们那里得酬报。那么这位神是需要酬报吗？但神身上竟有这一切，乃是侮辱。\par

在他之后，他们提出一位女神\Person{阿尔忒弥斯}，\Person{阿波罗}的姐妹；他们说她是一名女猎手，她自己常携带弓箭和箭矢，在山上漫游，带领猎犬猎取鹿或田野的野猪。但一个贞洁的处女独自在山上漫游或猎捕动物，这是可耻的。因此，\Person{阿尔忒弥斯}不可能是女神。\par

他们又说\Person{阿佛洛狄忒}确实是一位女神。她有时与他们的神同住，有时却又与人为邻。她曾以\Person{阿瑞斯}为情人，后又以\Person{塔慕兹}即\Person{阿多尼斯}为情人。\Person{阿佛洛狄忒}也曾为\Person{塔慕兹}之死哀哭切齿，他们说她还下到\Place{阴府}，为要从\Place{阴府}（\Person{哈得斯}）的女儿\Person{珀耳塞福涅}那里赎回\Person{阿多尼斯}。如果\Person{阿佛洛狄忒}是女神，却不能在情人死时帮助他，又怎能帮助别人呢？这是不能听的，一个神性竟会哭泣哀号并犯奸淫。\par

他们又说\Person{塔慕兹}是神。他，哼！是个猎人和奸夫。他们说他被野猪所伤而死，不能自救。如果他连自己都救不了，怎能顾念人类呢？但神是奸夫、猎人，或死于暴力，这是不可能的。\par

他们又说\Person{瑞亚}是他们众神的母亲。他们说，她曾有个情人\Person{阿提斯}，她喜欢堕落的男子。最后她为她的情人\Person{阿提斯}举哀悲恸。如果众神之母无法帮助她的情人、救他脱离死亡，她怎能帮助别人呢？所以，女神竟会悲恸哀哭并喜欢堕落的男子，这是可耻的。\par

他们又提出\Person{科瑞}，说她是女神，她被\Person{普路托}抢走，无法自救。如果她是女神且无法自救，她又怎能找到帮助别人的方法呢？因为一个被抢走的神是十分无能的。\par

王啊，这一切就是希腊人关于他们的神所提出的，他们捏造并宣告了这些事。由此，所有人受到冲动，去行各种亵渎和污秽的事；因而全地都败坏了。\par

12. 埃及人更甚，他们比地上一切民族都更卑劣愚蠢，自己也比一切人更迷失。因为野蛮人和希腊人的神灵（或宗教）尚不足以满足他们，他们竟又引入一些具有动物本性的东西，说它们是神，同样也把旱地上和水中的爬行动物当作神。他们还把一些植物和草木说成是神。他们比地上一切民族更被各种虚妄和污秽所败坏。因为从古时起，他们就敬拜\Person{伊西斯}，他们说她是一位女神，她的丈夫是她的兄弟\Person{奥西里斯}。当\Person{奥西里斯}被他的兄弟\Person{堤丰}杀死时，\Person{伊西斯}带着她的儿子\Person{荷鲁斯}逃往\Place{叙利亚}的\Place{比布鲁斯}，在那里住了一些时日，直到她的儿子长大。\Person{荷鲁斯}与他的叔叔\Person{堤丰}争战，杀了他。然后\Person{伊西斯}返回，带着她的儿子\Person{荷鲁斯}往来寻找她主\Person{奥西里斯}的尸体，为他的死痛哭哀号。如果\Person{伊西斯}是女神，却不能帮助她的兄弟和主\Person{奥西里斯}，她又怎能帮助别人呢？但神性会害怕、为安全而逃跑，或哭泣哀号，这是不可能的；否则那就极其悲惨。\par

他们又说\Person{奥西里斯}是一位有用的神。他被\Person{堤丰}杀死，无法自救。但众所周知，这不能用来断言神性。此外，他们说他的兄弟\Person{堤丰}是神，他杀了自己的兄弟，又被兄弟的儿子和新娘所杀，也不能自救。请问，一个不能救自己的人，怎能是神呢？\par

因此，既然埃及人比其余民族更为愚昧，这些以及此类神祇不能满足他们。不止如此，他们甚至将神的名号加在那些完全没有灵魂的动物身上。因为他们有的崇拜羊，有的崇拜牛犊；有的崇拜猪，有的崇拜鲥鱼；有的崇拜鳄鱼、鹰、鱼、朱鹭、秃鹫、雕和渡鸦。他们有的崇拜猫，有的崇拜大菱鲆，有的崇拜狗，有的崇拜蝰蛇，有的崇拜角蝰，还有的崇拜狮子；另有些崇拜大蒜、洋葱、荆棘，还有些崇拜老虎及其他诸如此类的事。可怜的人看不到所有这一切都算不得什么，尽管他们日日看见他们的神被人及同类吃掉、消耗；也看见有些被火化，有些死亡、腐烂而化为尘土，他们却没有注意到它们以各种方式灭尽。所以埃及人没有认识到，那些连自身得救都无能为力的东西，就不是神。如果它们在自救的事上尚且软弱，那么它们又有何能力帮助它们的敬拜者得救呢？埃及人陷入的错谬是何等的大啊——确实，比地面上任何民族的错谬都更大。\par

13. 然而，王啊，关于希腊人，真令人惊奇，他们在生活方式和理智上超越所有其他的民族，却如何误入歧途，追随死去的偶像和没有生命的像。他们明明看见他们的神在工匠手中被锯开、刨平、截短、砍削、烧焦，并受到种种装饰和改变。当它们变老，随时间消磨损毁，及至被熔化和碾成粉末时，我诧异他们怎能不觉察到它们不是神？至于那些连自己都无法拯救的，又怎能解救人的困苦呢？\par

但甚至他们中间的作家和哲学家也错误地宣称，神是为了尊崇全能的天主而造出的那类事物。他们犯错，试图将（它们）比作天主，而人从未在任何时候见过天主，也无法看到祂是什么样子。同样，他们（也犯错）在断言神性上会有匮乏之类的事；正如他们说祂接受祭品，需要燔祭、奠祭、人祭和庙宇。但天主并无所缺，这些事物对祂来说全无必要；显然人们在想象这些事上犯了错。\par

此外，他们的作家和哲学家声称并宣称他们所有神的本性是同一的。他们还没有领悟到天主我们的主，祂虽是一，却在万有之中。因此他们错了。因为如同人的身体虽有许多肢体，却一个肢体不畏惧另一个肢体，而是由于它是一个联合的身体，就全然与自身和谐一致；即使如此，天主在其本性上也是一。单一本质为祂独有，因为祂在其本性和本质上是统一的；祂不畏惧自己。那么，如果众神的本性是同一的，一位神就不应追赶、杀害或伤害另一位神。那么，如果有神被其他神追赶而受伤，有些被绑架，有些被雷电击死，很明显他们众神的本性不是同一的。由此可知，王啊，当他们将众神的本性归算并归入单一本性时，这是一个错误。那么，如果我们理应赞赏一位看得见却不能看的神，那么相信不可见且全见的本性，该是何等更值得称颂呢？而且，如果宜当赞许一位工匠的手艺，那么又该如何更宜当光荣那创造工匠的造物主呢？\par

看哪！当希腊人制定法律时，他们没有意识到，他们用自己的法律定了他们神祇的罪。因为如果他们的法律是正义的，他们的神就是不义的，因为他们触犯法律，互相残杀、行邪术、犯奸淫、抢劫偷盗、与男性同寝，并犯下其他种种恶行。因为如果他们的神做这一切正如所描述的那样，是正确的，那么希腊人的法律就是不义的，因为没有按照他们神祇的旨意来制定。而在那种情况下，全世界都误入歧途了。\par

关于他们神祇的叙述，有些是神话，有些是\OriginalTerm{自然诗}{\GreekText{φυσικαί}}，有些是颂歌和哀歌。颂歌和哀歌诚然是空洞的话语和噪音。但这些自然诗，即使如他们所说的被造出来，那些做这些事并遭受和忍受这些事的仍然不是神。那些神话是浅薄的故事，全无深度。\par

14. 王啊，现在让我们也来看看犹太人的历史，并看他们对天主有何意见。犹太人于是说，天主是唯一的，万物的造物主，全能；并且，除了这唯一的天主外，不应敬拜任何其他。在这点上，他们似乎比所有民族更接近真理，特别是在他们敬拜天主而非祂的造物。他们通过在他们中间盛行的仁爱来效法天主；因为他们怜悯穷人，释放被掳者，埋葬死者，并做这些在天主前蒙悦纳、在人前也蒙喜悦的事——这些（习俗）是他们从祖先领受的。\par

然而，他们也偏离了真知识。在他们的想象中，他们以为侍奉的是天主；但按他们的遵行方式，他们侍奉的是天使而非天主：——就如当他们庆祝安息日和月朔，无酵饼节期和大斋期的时候；还有禁食、割损和食物的洁净，然而这些事，他们并没有完全遵守。\par

15. 然而，基督徒啊，王啊，他们在四处奔走探索之后，寻得了真理；而且，我们从他们的经书得知，他们比起其余的民族，更接近真理和真正的知识。因为他们认识并信赖天主，即创造天地的造物主，万有都在祂内、由祂而有，没有别的神明与祂为伍；他们从祂领受了诫命，并将之铭刻于心中，怀着对来世的希望和期待而遵行。因此，他们不行奸淫，不犯邪淫，不作假见证，不侵吞抵押之物，不贪求不属于自己的东西。他们孝敬父母，善待近人；他们当审判者时，秉公判断。他们不敬拜按人形象造的偶像；凡自己不愿别人怎样对待自己的，他们也决不对别人那样做；对于祭偶像的食物，他们不吃，因为他们是洁净的。他们安抚压迫他们的人，使他们成为自己的朋友；他们善待仇敌；王啊，他们的妇女如同贞女般洁净，他们的女儿谦和端正；他们的男子远避一切不法结合和所有污秽，盼望在来世领受赏报。再者，他们中若有人蓄有奴仆或婢女或儿女，出于对他们的爱，便劝导他们成为基督徒；当他们这样做之后，便不分彼此地称他们为弟兄。他们不敬拜外邦的神明，他们行事为人温良而喜乐。在他们中间找不到虚谎；他们彼此相爱，对寡妇不失尊敬，并拯救孤儿脱离欺压者。有财产的，不求夸耀地分给那没有的。他们见到外乡人，便接进家中，喜乐地待之如同亲兄弟；因为他们在\OriginalTerm{精神}{spirit}内、在天主内称之为弟兄，不称血肉的弟兄。他们穷人中，若有一位离世，各人便按自己的能力去照顾他，细心办理安葬。他们若听说某人因他们的\OriginalTerm{默西亚}{Messiah}之名而被囚或受苦，众人便急切地照料他的急需，若有可能赎回他，就使他获得自由。他们中间若有贫穷匮乏的人，而他们没有多余的食物，便禁食两三天，以便补足穷人的匮乏。他们极其用心地遵守他们\OriginalTerm{默西亚}{Messiah}的诫命，公义而节制地生活，正如主他们的天主所吩咐的。每天早晨和每个时辰，他们都为天主向他们所施的慈爱而感谢赞美天主；也为他们的饮食而向祂感恩。他们中间的义人离世时，他们便欢喜并感谢天主；他们伴送遗体，仿佛他正从一处地方前往相近的另一处。他们中若有人生了孩子，就感谢天主；若这孩子竟在童年夭折，他们便更加感谢天主，如同为一位无罪的过客离世。再者，他们若见有人在不虔敬或罪恶中死去，就为他深深痛悔，忧伤如同为他正走向灭亡。\par

16. 王啊，基督徒的律法之命便是如此，他们的生活方式便是这样。他们作为认识天主的人，向祂祈求合乎祂赐予、也合乎他们领受的事物。他们就这样度尽一生。既然他们知道天主对他们的慈爱，看啊！为了他们的缘故，世上荣耀的事物便彰显出来。确实，他们在四处奔走寻求时，找到了真理；从我们所思考的来看，我们得知唯有他们接近真知。他们不在众人耳边宣扬自己所做的善事，而是谨慎不让人察觉；他们施舍隐藏起来，如同人寻得宝藏而藏匿一般。他们努力追求公义，如同那些期待仰见他们的\OriginalTerm{默西亚}{Messiah}、并从祂那里以极大的荣耀领受关于他们的应许的人。至于他们的言语和诫命，王啊，以及他们在敬拜中的荣耀，和他们按各人的工作希望于来世获得赏报的期待，你从他们的经书便可得知。我们简要地向陛下陈明了基督徒的行为和真理，这已足够。因为他们的教义，凡深入查考并反思的人，必觉其伟大而奇妙。确实，这是一个新民族，在他们中间有某种\OriginalTerm{神圣的混合}{a divine admixture}存在。\par

那么，你拿他们的经书来读吧，看啊！你会发现，我并非凭自己的权威提出这些事，也不是作为他们的辩护者才这样说；而是因为我读了他们的经书，完全确信这些事以及将来之事。为此，我不得不向那些关心真理、寻求来世的人宣告真理。我毫不怀疑，大地因基督徒的祈求而得以存立。然而，其余的民族却误入歧途，使人迷误，他们沉迷于世上的元素，因为他们的神识不能超越于此。他们如同在黑暗中摸索，因为不愿认识真理；又像醉酒之人，东倒西歪，彼此推撞，跌倒在地。\par

17. 王啊，我就说到这里；至于其余的，正如上文所说，在他们的其它经书中，记载着一些难以言说、难以叙述的事——这些事不仅在言语中说出，也以行为彰显出来。\par

王啊，希腊人行可耻的事，与男性、与母亲、与姐妹、与女儿交合，却反过来将他们可怕的不洁归咎于基督徒。然而，基督徒是公义良善的，真理摆在他们眼前，他们的心灵坚忍；因此，尽管他们知道这些人的错谬，并受其逼迫，他们还是忍受担当；他们多半怜悯这些人，视之为缺乏知识的人。并且，他们也为这些人祈祷，愿他们悔改错谬；若有人悔改了，他就在基督徒面前为自己的所作所为感到羞愧；他向天主\OriginalTerm{告解}{confession}说：\Quote{我是在无知中做了这些事。}他洁净自己的心，他的罪就蒙赦免，因为他是在从前无知的时候犯下的，那时他常亵渎并毁谤基督徒的真知识。诚然，基督徒的族类比地上所有的人更为有福。\par

从此，让那些说虚妄话、骚扰基督徒的舌头闭口不言吧；今后让他们说真理。因为对他们而言，崇拜真天主远比崇拜无意义的声音要紧得多。的确，基督徒口中所说的，都是出于天主；他们的教义是光明之门。因此，凡不认识天主的人，都要来靠近它；他们必领受不朽的言语，那些言语出自万古，来自永恒。如此，他们便能面对那通过\OriginalTerm{耶稣默西亚}{Jesus the Messiah}预定要临到全人类的可畏审判。\par

\WorkTitle{哲学家\Person{阿里斯蒂德}的护教篇}终。\par

\SourceRecord{Translated by D.M. Kay.；Ante-Nicene Fathers；Vol. 9.；Edited by Allan Menzies.；Buffalo, NY: Christian Literature Publishing Co.,；1896.；<http://www.newadvent.org/fathers/1012.htm>.}
